\chapter*{Введение}
\addcontentsline{toc}{chapter}{Введение}

JWT токены получили широкое распространение как механизм авторизации в распределённых
микросервисных системах~\cite{RFC6749}. % Не согласен с цитированием, почему тут oauth&
Надёжность схемы с асимметричными ключами целиком определяется защищённостью приватного ключа подписи:
его компрометация позволяет злоумышленнику выдавать произвольные токены и получать
несанкционированный доступ к ресурсам всех потребителей системы~\cite{Arora2024}.
В распределённых платформах широко распространена практика, при которой единственный ключ подписи
доступен всем микросервисам одновременно и не ротируется, — что существенно расширяет
поверхность атаки~\cite{Matias2023}. % Вряд ли в этм источнике что-то подобное описанл
При этом реализация плановой ротации ключей в высоконагруженных системах с большим числом сервисов
сопряжена с необходимостью обеспечить непрерывность обработки запросов в процессе смены ключей
и согласованность состояния всех компонентов.

\textbf{Объектом исследования} является процесс межсервисной авторизации в высоконагруженной
микросервисной платформе ООО~<<Майндбокс>>, функционирующей на базе Kubernetes.

\textbf{Предметом исследования} выступает система автоматической бесшовной ротации
криптографических ключей авторизации указанной платформы.

\textbf{Целью работы} является проектирование и реализация системы автоматической бесшовной
ротации криптографических ключей авторизации в высоконагруженной микросервисной платформе.

Для достижения поставленной цели необходимо решить следующие \textbf{задачи}:
\begin{enumerate}
  \item Проанализировать текущую систему межсервисной авторизации и выявить её ограничения.
  \item Исследовать существующие решения для бесшовной ротации криптографических ключей.
  \item Спроектировать архитектуру системы ротации ключей с поддержкой межконтурного взаимодействия.
  \item Разработать и протестировать компоненты системы.
  \item Мигрировать микросервисы платформы на новую систему авторизации.
  \item Провести апробацию системы: запустить автоматическую ротацию ключей по~расписанию.
\end{enumerate}

% В гипотезе не нужны такие детали.
% Здесь уместнее говорить языком метрик чего мы достигнем, какую проблему закроем
В основу работы положена \textbf{гипотеза} о том, что разработка собственного решения
является единственным приемлемым способом обеспечить бесшовную автоматическую ротацию ключей
в рамках существующей архитектуры платформы.
Анализ рыночных альтернатив показал их непригодность: HashiCorp Vault обладает ограниченной
масштабируемостью и не поддерживает необходимое межконтурное взаимодействие;
Keycloak не имеет встроенных механизмов ротации ключей подписи, а~разработка расширения
на~Java несовместима с~технологическим стеком платформы (C\#);
переход на~Service Mesh с~взаимной TLS-аутентификацией потребовал бы полной переработки
сетевой инфраструктуры, что неприемлемо с~точки зрения оперативного закрытия уязвимости.
Собственное решение позволяет не~только устранить выявленную уязвимость, но~и~получить
ряд дополнительных преимуществ: централизованное управление ключами, изоляцию контуров
развёртывания и~ограниченный срок жизни токенов.

\textbf{Практическая значимость} работы определяется тем, что разработанная система
внедрена в промышленную эксплуатацию в платформе, обрабатывающей свыше четырёх миллионов запросов в минуту.
Реализованное решение закрывает критическую уязвимость: доступ к приватному ключу подписи JWT
теперь ограничен единственным сервисом авторизации, тогда как ранее он был доступен всем 120 микросервисам платформы.
Ключи ротируются автоматически по расписанию без прерывания работы сервисов
при сохранении сквозного SLA на генерацию токена не более 30 мс по 95 процентилю.
Дополнительно обеспечены изоляция контуров развёртывания, централизованное управление
авторизационными клиентами через IaC-репозиторий и ограниченный срок жизни токенов,
что в совокупности существенно снижает риски при компрометации отдельных компонентов системы.
